\chapter{Preliminaries}
\label{chap:preliminaries}

\section{Definitions}
\label{sec:definitions}

This section establishes the terminology used throughout the thesis. The definitions are operational rather than purely theoretical: they specify how each concept is understood within the analysis of arrival management at Blaise Diagne International Airport. Establishing this terminology is necessary because the analytical, forecasting, simulation, and decision-support components of the study examine the same operational system from different perspectives. A consistent vocabulary, therefore, provides the basis for interpreting the methodology and results presented in the subsequent chapters.

\subsection{Operational terms}
\label{subsec:operational-terms}

Airport arrival operations refer to the coordinated processes through which arriving passengers, their baggage, and the associated operational information move from aircraft arrival to completion of the terminal arrival journey. Within this thesis, the arrival process comprises passenger generation following a commercial flight arrival, immigration processing, baggage reclaim, customs inspection, and passenger exit. These activities are supported by airport personnel, immigration and customs officers, airlines, ground-handling teams, baggage-handling personnel, security staff, and other operational stakeholders.

Arrival operations are treated as an interconnected system rather than as a set of independent service points. Passengers generally move through the processes in sequence, while baggage follows a related handling and reclaim process. Pressure arising at one stage can therefore affect the performance of subsequent stages. A delay at immigration may change when passengers reach baggage reclaim; slow baggage delivery may increase the number of passengers waiting in the reclaim area; and insufficient customs-processing capacity may cause queues to extend into earlier parts of the arrival system. The operational performance of the entire arrival journey consequently depends on both the capacity of individual processes and their coordination.

An arrival bank is a period in which several commercial flights arrive within a relatively short time window. Arrival banks concentrate passenger and baggage demand because passengers from multiple flights enter the arrival system at approximately the same time. They are operationally significant at AIBD because the resulting demand may simultaneously pressure immigration counters, baggage-reclaim facilities, customs checkpoints, and terminal exits.

The term arrival flow denotes the movement of passengers through the successive stages of the arrival system. It includes both the volume of passengers entering each process and the timing of their movement between processes. Arrival flow is not assumed to be uniform. It varies according to the number and timing of arriving flights, the estimated passenger volume associated with those flights, baggage demand, processing conditions, and the resources available at each service point.

Passenger demand is the estimated number of passengers entering the arrival system over a defined period. In this study, passenger demand is derived from commercial flight arrival records and aggregated into operationally relevant time intervals. It represents the workload imposed on arrival-side services, particularly immigration, baggage reclaim, customs, and passenger-exit processes.

Passenger pressure describes the operational intensity created by the number and concentration of arriving passengers during a given period. It is broader than passenger volume alone. Passenger pressure reflects the extent to which the timing and concentration of demand may challenge available processing capacity, leading to queues, longer processing times, or congestion. A period with a high concentration of arriving passengers is therefore operationally more demanding than a period in which a comparable volume is distributed more evenly.

Baggage demand refers to the estimated baggage load associated with arriving commercial passengers. It represents the quantity of baggage that must be transferred from arriving aircraft, handled through the airport's baggage processes, and delivered to reclaim facilities. In the thesis, baggage demand is analyzed separately from passenger demand because the two are related but operationally distinct. Passenger processing and baggage delivery may be subject to different capacities, service conditions, and sources of delay.

Baggage pressure is the operational intensity placed on baggage-handling and reclaim processes by the volume and concentration of baggage demand. High baggage pressure may occur when several flights require baggage-handling capacity simultaneously. Its consequences may include increased baggage queues, slower delivery, longer passenger waiting times in the reclaim area, and additional pressure on the wider arrival process.

The distinction between passenger pressure and baggage pressure is central to the thesis. Passenger volume indicates the expected demand on immigration, customs, and other passenger-processing activities, whereas baggage demand indicates the expected workload for baggage-handling and reclaim operations. Although both pressures originate from arriving flights, they need not produce identical operational effects. Their separate treatment permits a more precise assessment of arrival conditions and supports the development of process-specific recommendations.

Operational capacity refers to the amount of passenger or baggage demand that an arrival process can handle under a given resource configuration. Capacity is determined by the resources assigned to a process and by the rate at which those resources can complete their activities. In the context of this thesis, relevant capacity elements include immigration counters, baggage-processing and reclaim capacity, customs-processing resources, and the staff required to operate or coordinate these services.

Demand-capacity imbalance occurs when the passenger or baggage demand entering a process exceeds the capacity of the available resources within the same period. Such an imbalance does not necessarily imply permanently insufficient capacity. It may instead result from the temporary concentration of flights during an arrival bank. Where demand exceeds processing capability, unfinished work accumulates and forms a queue.

A queue is the accumulation of passengers or baggage awaiting service because immediate processing is unavailable. The study considers queues at immigration, baggage reclaim, and customs. Queue size indicates the amount of unmet demand at a particular process. Both arrival intensity and service capacity influence queue development: queues grow when arrivals occur faster than they can be processed and decline when processing capacity exceeds incoming demand.

Congestion refers to an operational condition in which concentrated passenger or baggage demand results in sustained queues, increased system times, pressure on processing resources, or interference among connected arrival processes. Congestion is therefore treated as a system condition rather than as the mere presence of a short queue. It becomes operationally relevant when demand-capacity imbalance affects the efficiency, predictability, or quality of the arrival process.

A congestion period is a time interval during which passenger pressure, baggage pressure, queue accumulation, or processing delays indicate an elevated risk of operational bottlenecks. Congestion periods may be associated with arrival banks, limited service capacity, delayed processing, or a combination of these conditions. Identifying such periods in advance is a primary purpose of the predictive framework developed in this thesis.

A bottleneck is an arrival-process stage whose available processing capability restricts the performance of the wider system. Immigration, baggage reclaim, and customs can each become bottlenecks when demand temporarily exceeds the capacity available at that stage. Because the arrival processes are interconnected, a bottleneck may affect not only the passengers immediately waiting for service but also the timing and workload of downstream operations.

Passenger system time is the total time a passenger spends within the modeled arrival system, from entry into the arrival process until completion of customs processing and exit. It captures the combined effects of waiting and processing across the arrival journey. Passenger system time is therefore used as an overall indicator of arrival-process performance rather than as a measure of a single service point.

Processing time refers to the time required to complete a specific operational activity, such as immigration control or customs inspection. It is distinguished from waiting time, during which a passenger is ready for processing but cannot yet access the required service. Extended total system time may result from long processing or waiting times, or both.

Passenger throughput denotes the number of passengers who complete the modeled arrival process within a defined period. Throughput indicates the effective output of the arrival system under a particular demand level and resource configuration. It is used to assess whether operational changes increase the system's ability to process arriving passengers.

Operational resilience refers to the capacity of the arrival system to maintain or recover acceptable performance under increased demand, peak-arrival conditions, or operational stress. Within the simulation component of the study, resilience is assessed by examining how queues, system times, throughput, and congestion severity respond under different scenarios. A more resilient configuration can absorb or manage pressure without a disproportionate deterioration in arrival performance.

Operational interpretability is the extent to which an analytical or predictive output can be understood in terms of airport processes and translated into a relevant management decision. A forecast is operationally interpretable when it indicates not only that pressure is expected, but also when it is likely to occur, which part of the arrival process may be affected, and what type of resource or intervention may require consideration.

The term Teranga is used in the thesis to refer to the human-centered hospitality associated with Senegal \parencite{demir2025speaker}. Within the AIBD context, Teranga connects the airport's operational objectives with its role as an international gateway. Efficient arrival management is therefore not merely a question of processing speed. It also supports a smooth, welcoming, and predictable passenger experience consistent with the values the airport seeks to represent.

\subsection{Analytical and decision-support terms}
\label{subsec:analytical-decision-support-terms}

Operational data are records describing activities, events, or conditions within the airport system. The operational data used in this thesis consist primarily of commercial flight arrival records and the temporal and operational variables required to estimate passenger and baggage demand and arrival pressure.

The study is based on a validated master dataset containing 22,270 records of commercial passenger arrivals. The dataset provides the analytical foundation for examining temporal arrival patterns, constructing passenger and baggage indicators, identifying pressure periods, developing forecasts, preparing simulation inputs, and producing decision-support outputs. The use of a validated master dataset ensures that the forecasting, simulation, and dashboard components are based on a consistent set of arrival records.

FR24-like data refers to the commercial flight arrival data used in the study, organized in a format comparable to flight-tracking arrival records. The term refers to the operational structure of the dataset without altering the nature of the underlying records. The data describe commercial arrivals and provide the flight-level and temporal information from which the study derives passenger, baggage, and pressure indicators.

A temporal variable is a variable that represents the time of an arrival event. Temporal variables support the analysis of variation across arrival hours and other available time-related groupings. They are used to determine whether passenger and baggage demand follow recurring patterns and to identify periods associated with increased pressure.

An operational variable is an available characteristic of a commercial arrival relevant to passenger demand, baggage demand, or the behavior of the arrival system. Temporal and operational variables are considered together because pressure is influenced not only by the number of arriving flights but also by their characteristics and concentration over time.

A pressure indicator is a derived operational measure used to represent passenger demand, baggage demand, or congestion risk in a form suitable for analysis, forecasting, simulation, and reporting. Pressure indicators convert individual arrival records into measures that can be interpreted at the level of an operational time interval. They create a common analytical link between the source data and the later stages of the framework.

A pressure classification is the assignment of an operational pressure state to an analyzed time interval based on the indicators developed in the study. Pressure classifications support the distinction between routine conditions and periods requiring increased operational attention. They are used to improve the interpretability of model outputs and to support dashboard-based risk communication.

Predictive operational analytics refers to the application of analytical and forecasting methods to operational airport data to anticipate future arrival conditions. In this thesis, predictive operational analytics is used to estimate passenger pressure, baggage demand, and congestion periods before they develop into bottlenecks.

Predictive operational intelligence is the broader framework that integrates operational data, pressure indicators, predictive models, simulation results, key performance indicators, and decision-support outputs. Predictive operational intelligence extends beyond the production of forecasts. It connects prediction with operational evaluation and management action. Its purpose is to move arrival management from retrospective monitoring and reactive intervention towards advanced preparation and proactive decision-making.

A forecast is an estimate of a future operational value or state based on the temporal and operational information available in the arrival data. The forecasts developed in the thesis concern arrival-passenger pressure, baggage demand, and congestion periods. Their value is assessed not only through forecasting accuracy but also through their relevance to resource planning and operational intervention.

The forecasting horizon is the interval between the time at which a prediction is produced and the future period to which it refers. In this thesis, the horizon is short, because a forecast is operationally useful only if it arrives early enough for staffing, baggage-capacity, or queue-management action to be taken before the pressure materializes. Forecasting accuracy is correspondingly assessed against that operational purpose rather than against statistical error alone.

Discrete-event simulation represents the arrival system as a sequence of events that occur at particular points in simulated time. Relevant events include passenger entry into the system, arrival at a service point, entry into a queue, commencement of processing, completion of processing, movement to a subsequent stage, and exit from the system. The approach is appropriate to the thesis because passenger movements and service completions occur as distinct operational events, and queues develop through their interaction.

A key performance indicator, or KPI, is a defined measure used to evaluate the behavior and performance of the simulated arrival system; the indicators used in this thesis are specified in Chapter 4, Section 4.5. Congestion severity is the KPI representing the extent to which arrival pressure affects the operation of the modeled system.

A decision-support system, or DSS, is a structured environment that integrates data and analytical outputs to assist managerial judgment. The DSS developed in the thesis does not replace airport managers or operational stakeholders. It organizes forecasts, pressure indicators, simulation results, KPIs, risk information, scenario comparisons, and recommendations so that decision-makers can evaluate operational conditions and possible responses.

The Power BI Decision Support System is the manager-facing component of the predictive operational intelligence framework. It translates the outputs of the operational analysis, forecasting models, and AnyLogic simulation into interactive dashboard views. These views support operational monitoring, congestion-risk assessment, scenario comparison, dynamic-intervention planning, executive decision-making, and the interpretation of factors associated with operational pressure.

A manager-facing insight is an analytical finding expressed in a form that corresponds to a practical airport decision. For example, a predicted pressure period becomes manager-facing when it is linked to the timing of the pressure, the process likely to be affected, the expected operational consequences, and the type of staffing, baggage capacity, or queue management response that may be required.

An actionable recommendation is a decision-support output that identifies an operational response connected to the predicted or simulated condition. Actionability is essential because a model may produce technically valid information that remains of limited practical value if airport stakeholders cannot determine how to use it.

\section{Conceptual Foundations}
\label{sec:conceptual-foundations}

The thesis is founded on the view that an airport arrival terminal functions as an interconnected operational system. Its performance arises from the interaction among scheduled arrivals, passenger demand, baggage demand, service resources, processing times, queues, and decisions made by multiple stakeholders. The study, therefore, does not treat congestion as an isolated event at a single checkpoint. It examines how demand enters the system, how that demand is distributed across connected processes, how capacity limitations cause waiting to accumulate, and how analytical information can support intervention before operational performance deteriorates.

The first conceptual building block is the distinction between monitoring and prediction. Monitoring describes current or recently observed operational conditions. At AIBD, real-time feedback mechanisms and operational alerts can help staff identify existing passenger service problems. Such information is important for immediate response, but it is inherently reactive when the problem has already developed.

Prediction has a different function. It estimates the passenger, baggage, or congestion pressure expected in a future operational period. This distinction is central to the research problem. The thesis does not seek to replace real-time monitoring; it seeks to complement it by providing advanced information. The integration of monitoring and forecasting enables identifying a developing pressure period before queues become severe and the passenger experience is adversely affected.

The second building block is the relationship between demand and capacity. Passenger and baggage demand originate from commercial flight arrivals. When flights are sufficiently separated, the arrival system may process the associated demand without substantial accumulation. When several flights converge, the same resources must serve a larger number of passengers and handle a larger baggage load within a shorter interval.

Operational pressure arises from this temporal concentration. The total daily passenger volume alone is therefore insufficient for understanding congestion. Two days with comparable passenger totals may create different operational conditions if one contains concentrated arrival banks and the other distributes demand more evenly. For arrival management, the timing and intensity of demand are consequently as important as its overall magnitude.

Capacity is also process-specific. Immigration capacity is not interchangeable with baggage-reclaim or customs capacity. Increasing the number of available immigration counters may reduce immigration queues, but it may also cause passengers to reach baggage reclaim more quickly. If baggage capacity remains unchanged, pressure may shift between stages. Similarly, improving baggage delivery without sufficient customs-processing capacity may increase downstream queues. Operational improvement must therefore be assessed at the system level rather than inferred from the performance of a single process.

The third building block concerns queue formation and propagation. A queue forms when passengers or baggage arrive faster than the corresponding process can complete service. The queue represents accumulated demand that has not yet been processed. Its size at any point in time reflects the earlier interaction between arrivals and service completions.

Queue size alone, however, does not describe the full operational effect. A queue may be short-lived and absorbed quickly, or it may persist and affect connected processes. Persistent immigration queues increase passenger waiting time and delay movement towards baggage reclaim. Baggage delays keep passengers within the reclaim area and may reduce the predictability of customs demand. Customs queues delay final exit and may extend into upstream spaces. Congestion can therefore propagate through the arrival journey even when its initial cause is associated with one service point.

This system interaction explains the need for several KPIs. Passenger system time measures the combined effect of waiting and processing throughout the entire arrival journey. Process-specific queue sizes identify where unmet demand is accumulating. Throughput measures how many passengers complete the system. Congestion severity describes the overall intensity of operational pressure, while resilience indicates how effectively the system maintains or restores performance under stress. No single KPI is sufficient to evaluate the complete operational condition.

The fourth conceptual building block is short-horizon forecasting. The forecasting task is operational rather than long-range or strategic. AIBD requires information about approaching passenger and baggage pressure early enough to adjust resources or prepare interventions. The relevant horizon is therefore close to the period of operation, while still preceding the anticipated bottleneck.

Short-horizon forecasting establishes a practical link between flight arrival data and airport management. The arrival records contain temporal and operational information that can be used to estimate future passenger and baggage pressure. Machine-learning models are then used to identify relationships between those inputs and the operational outcomes of interest.

The conceptual sequence is not limited to producing a numerical prediction. A forecast must pass through several stages before it can support airport operations. First, the source data must be validated and structured. Second, passenger and baggage demand must be represented through operationally meaningful indicators. Third, the relevant temporal and operational features must be prepared for modeling. Fourth, the forecasting model must produce an estimate or classification. Finally, that output must be interpreted in relation to airport processes, capacity, and possible management action.

This sequence distinguishes predictive accuracy from operational utility. A model may produce small statistical errors but remain difficult to use if its outputs do not indicate when attention is required or how the predicted pressure relates to an operational process. Conversely, an interpretable pressure classification may support decisions more directly by communicating the timing and significance of the expected condition. The thesis, therefore, considers forecasting accuracy and operational interpretability together.

The fifth building block is simulation-based operational evaluation. Historical analysis can reveal recurring patterns, and predictive models can estimate future demand, but neither alone establishes how the arrival system will behave under a particular resource configuration. Simulation provides this additional layer by representing the movement of passengers through immigration, baggage reclaim, customs, and exit while recording waiting, processing, queue accumulation, and completion.

The AnyLogic model serves as an analytical bridge between forecast demand and operational consequences. Predicted or scenario-specific pressure is entered into the model as a demand condition. The model then represents how that demand interacts with processing capacities and service stages over simulated time. The resulting KPIs indicate whether the configuration absorbs the pressure or develops bottlenecks.

The scenario structure supports different forms of operational reasoning. The baseline scenario establishes a reference point. The capacity-optimization scenario evaluates changes to available resources. The peak-stress scenario tests the system under intensified demand. The dynamic-intervention scenario examines responses associated with changing pressure. The integrated-optimization scenario assesses the combined effect of the improvements considered in the framework.

This approach does not assume that one intervention is effective under all conditions. An increase in capacity may produce substantial benefits during a peak period but have limited value under routine demand. An intervention applied to one process may reduce one queue while transferring pressure to another. Scenario comparison is therefore required to evaluate both the direct and system-wide consequences of operational changes.

The simulation also supports the concept of operational resilience. Resilience cannot be determined solely from performance under routine conditions. It becomes visible when the system is exposed to stress, and its response is compared across configurations. If queues and system times increase sharply during peak demand, the configuration has a limited ability to absorb the pressure. If an intervention constrains that deterioration and supports recovery, it contributes to a more resilient arrival process.

The sixth conceptual building block is decision support. Analytical models do not make operational decisions independently. Airport managers and partner organizations remain responsible for evaluating conditions, coordinating resources, and selecting appropriate actions. The purpose of the DSS is to make the relevant evidence accessible and interpretable at the point of decision.

The decision-support logic of the thesis comprises four interconnected elements. The first is awareness: the system must communicate the current or expected level of passenger and baggage pressure. The second is diagnosis: it must indicate which process or operational factor is associated with the condition. The third is evaluation: managers must be able to compare scenarios and examine the expected consequences of alternative configurations. The fourth is action: the analytical output must be linked to staffing, baggage capacity, queue management, or related resource-planning decisions.

Power BI provides the environment through which these elements are brought together. The dashboards consolidate forecasting outputs, simulation metrics, congestion information, scenario comparisons, and recommendations. The purpose is not simply to display a large number of measures. The dashboard must organize them according to operational questions: when pressure is expected, how severe it may become, where it is likely to occur, how alternative scenarios perform, and which management response is supported by the available evidence.

Interpretability is particularly important in a multi-agency airport environment. Immigration officers, baggage handlers, airline operations personnel, and airport managers each have different responsibilities. A technical prediction that is meaningful to a data analyst may not automatically correspond to the information required by each operational stakeholder. The DSS therefore translates analytical outputs into common pressure indicators, KPIs, and scenario-based findings that support coordination across roles.

The complete conceptual framework follows the sequence established in \cref{chap:introduction}: explain, predict, and support decisions. The explanatory stage identifies temporal and operational patterns in passenger and baggage demand. The predictive stage estimates future pressure and congestion periods. The decision-support stage uses simulation and dashboards to evaluate the operational meaning of those forecasts and present the results in an actionable form.

These stages are interdependent. Explanation without prediction remains retrospective. Prediction without simulation provides limited information about operational consequences. Simulation without decision support may remain inaccessible to managers. A dashboard without validated analytical inputs may communicate information clearly but lack a reliable evidence base. The thesis, therefore, integrates all four functions (operational analysis, forecasting, simulation, and decision support) within a single arrival-centric framework.

\section{Operational Context}
\label{sec:operational-context}

Blaise Diagne International Airport is the operational setting of the study. Since opening in 2017, the airport has developed as a strategic gateway for Senegal and the wider West African region, connecting Africa with Europe and the Americas. Its passenger volume increased from 2.6 million in 2022 to 2.94 million in 2024. The airport's stated development objectives include annual passenger volumes of five million by 2025 and ten million by 2035 \parencite{demir2025speaker}. This growth strengthens AIBD's regional role but also increases the operational demands on passenger processing, baggage handling, customs, and supporting services.

The convergence of several flows characterizes the airport's arrival environment. Aircraft arrive according to commercial schedules. Passengers disembark and enter immigration processing. Their baggage is transferred from the aircraft into the baggage-handling and reclaim system. Following immigration, passengers proceed to baggage reclaim, continue through customs, and leave the terminal. Each stage depends on the timely availability of staff, processing resources, operational information, and coordination between organizations.

The modeled passenger flow follows the sequence:

\begin{center}
\textbf{Commercial flight arrival $\rightarrow$ passenger generation $\rightarrow$ immigration processing $\rightarrow$ baggage reclaim $\rightarrow$ customs inspection $\rightarrow$ passenger exit.}
\end{center}

This sequence provides the operational structure for the thesis's analytical and simulation components. It does not imply that all passengers experience identical processing times or that the flow remains constant. The purpose is to represent the principal stages through which arrival demand moves and at which operational pressure may emerge.

Commercial arrivals provide the initial source of demand. Each flight contributes an estimated passenger volume and an associated baggage load. When flights arrive at separated times, demand is distributed across the available processing period. When several flights converge, passengers and baggage enter their respective processes in concentrated waves. These waves create arrival banks and increase the risk that available capacity will be temporarily insufficient.

Immigration is the first principal passenger-processing stage represented in the arrival journey. Its workload depends on the timing and estimated passenger volume of arriving flights. Immigration counters can become a bottleneck when the number of passengers requiring processing exceeds the capacity of available counters and officers. Queue buildup at this stage increases wait times and delays passenger movement towards baggage reclaim.

Immigration performance is also relevant to operational coordination. The rate at which immigration processes passengers influences when those passengers reach subsequent stages. A slower immigration process may temporarily limit pressure at baggage reclaim, but increase waiting at the border-control stage. A faster immigration process may reduce one queue while increasing the rate at which passengers enter the reclaim area. The evaluation of immigration capacity must therefore consider its effect on the complete arrival flow.

Baggage reclaim constitutes a related but separate operational stream. While passengers pass through immigration, baggage is unloaded, transferred, sorted, and delivered to reclaim facilities. Baggage demand is linked to commercial arrivals, but its processing time and capacity constraints differ from those of passenger control. The alignment between passenger release from immigration and baggage availability affects the time passengers spend in the reclaim area.

When several flights require baggage-processing capacity during the same period, baggage queues and delivery delays may develop. Such conditions affect passengers directly through extended wait times and affect airlines and ground-handling teams through baggage performance requirements. Baggage-reclaim pressure may also contribute to crowding and may alter the timing of demand reaching customs.

Customs is the final principal processing stage before passenger exit. Demand at customs reflects the flow of passengers completing the preceding stages. Customs queues may therefore arise from both the volume of passengers and the rate at which immigration and baggage processes release them. Insufficient customs-processing capacity can extend total system time even when earlier stages perform efficiently.

The operational system is thus sequential but not independent. Each process has its own resources and constraints, while its output becomes input to another part of the system. This interdependence explains why arrival management cannot be reduced to the separate monitoring of immigration, baggage, and customs. An intervention must be assessed based on its effect on the entire passenger journey.

The operational gap is therefore not an absence of all information. It is the limited availability of integrated short-horizon forecasts and manager-facing mechanisms that connect predicted pressure with resource decisions. Without such support, managers may recognize congestion only after queues have formed, requiring reactive adjustments under time pressure. Late intervention can lead to inconsistent service, inefficient staff deployment, and increased stress for employees and passengers.

\subsection{Stakeholders, constraints, and operational pressures}
\label{subsec:stakeholders-constraints-operational-pressures}

Arrival management at AIBD involves several stakeholder groups whose responsibilities are distinct but operationally interdependent. Immigration officers process passengers at the border-control stage and need advance information on passenger peaks to support counter allocation; customs personnel manage the final control stage, where workload depends on the flow released by immigration and baggage reclaim; airline operations and ground-handling teams depend on arrival processing and baggage delivery for schedule and service performance; and passengers are themselves not a uniform group, since citizens and tourists seek a smooth and welcoming arrival while business travellers require predictable processing times. These needs are connected: a passenger-centered outcome depends on each group's ability to work from a shared understanding of when pressure will occur.

Several constraints shape this operational environment. The first is traffic growth, which raises demand across every arrival stage. The second is the temporal concentration of that demand: resources sufficient for an average load may be inadequate when several flights converge. The third is process interdependence, since an adjustment at one stage changes the demand reaching the next. The fourth is multi-agency coordination, because immigration, customs, airport management, airlines, and handling teams operate under separate lines of authority. The fifth is the limited availability of predictive, manager-facing intelligence, with existing feedback being largely retrospective. The sixth concerns the balance between technological modernization and human-centered service, since AIBD's objective is to combine the two rather than trade one against the other.

Within these constraints, the study treats arrival pressure as a multidimensional operational issue. It affects service efficiency through queues and system times; resource use through staffing and capacity requirements; passenger experience through waiting and uncertainty; airline performance through baggage delivery and customer satisfaction; employee conditions through workload and reactive intervention; and national image through the quality of the airport's international arrival experience.

The predictive operational intelligence framework responds to this environment by connecting validated commercial arrival records with passenger and baggage indicators, short-horizon forecasting, AnyLogic simulation, operational KPIs, and Power BI decision support. The framework is designed to identify when pressure is likely to occur, examine how the arrival system may respond, compare possible operational configurations, and communicate the resulting evidence to managers and partner organizations.

This operational context also determines the scope of the subsequent literature review. The relevant body of research must be considered in relation to the specific requirements established in this chapter: interconnected airport arrival processes, temporary demand-capacity imbalance, passenger and baggage pressure, short-horizon prediction, queue and congestion behavior, simulation-based scenario evaluation, and manager-facing decision support. Chapter 3, therefore, positions the thesis within the existing work on airport operational analytics while examining the gaps it addresses through its arrival-centric and integrated framework.